%% Supplementary material for submission to JRSSB.
%% This file is standalone and uses the same OUP/RSS authoring template
%% as the main manuscript.

\documentclass[numsec,webpdf,modern,medium,namedate]{oup-authoring-template}

\onecolumn

\usepackage[doublespacing]{setspace}
\usepackage[fontsize=12pt]{fontsize}

% Theorem environments
\theoremstyle{thmstyleone}
\newtheorem{theorem}{Theorem}
\newtheorem{assumption}[theorem]{Assumption}
\newtheorem{lemma}[theorem]{Lemma}
\newtheorem{corollary}[theorem]{Corollary}
\theoremstyle{thmstyletwo}
\newtheorem{remark}[theorem]{Remark}

% Notation used in the supplement
\newcommand{\R}{\mathbb{R}}
\newcommand{\E}{\mathbb{E}}
\newcommand{\Prob}{\mathbb{P}}
\newcommand{\dd}{\mathrm{d}}
\newcommand{\Leb}{\mathrm{Leb}}

% Prefix every numbered object with S, to distinguish the supplement
\makeatletter
\def\@S#1{S\@arabic#1}
\makeatother
% from the main article.
\renewcommand{\thesection}{S.\arabic{section}}
\renewcommand{\thesubsection}{\thesection.\arabic{subsection}}
\renewcommand{\thesubsubsection}{\thesubsection.\arabic{subsubsection}}
\renewcommand{\theequation}{S.\arabic{equation}}
\renewcommand{\thetheorem}{S.\arabic{theorem}}
\renewcommand{\thefigure}{S.\arabic{figure}}
\renewcommand{\thetable}{S.\arabic{table}}
\renewcommand{\thefootnote}{S\arabic{footnote}}

\begin{document}

\journaltitle{Journal of the Royal Statistical Society Series B: Statistical Methodology}
\DOI{DOI HERE}
\copyrightyear{XXXX}
\pubyear{XXXX}
\access{Advance Access Publication Date: Day Month Year}
\appnotes{Supplementary material}

\firstpage{1}
\pagenumbering{S}

\title[Supplementary Material]{Supplementary Material for
``Bayes-assisted anytime-valid $t$-tests for Gaussian means with
unknown variance''}

\author[$\ast$]{Debolina Paul}
\author[]{Patrick Rebeschini}
\author[]{Fran\c cois Caron}
\authormark{Paul et al.}

\address[1]{\orgdiv{Department of Statistics},
  \orgname{University of Oxford},
  \orgaddress{\street{24--29 St Giles}, \postcode{OX1 3LB},
  \state{Oxford}, \country{UK}}}

\corresp[$\ast$]{Address for correspondence. Name, institution, city,
postcode, country. \href{mailto:email@example.com}{email@example.com}}

\abstract{This supplementary material extends the prior-misspecification
result in the main article to prior tails containing an exponential
factor.  After suitable normalization, a prior shifted arbitrarily far
from a fixed data-generating effect converges locally to an exponential
tilt of Lebesgue measure.  Extended Ville calibration consequently
converges to a directional, exponentially tilted benchmark.}

\keywords{Anytime-valid inference; extended Ville inequality;
prior misspecification; regular variation; sequential testing}

\maketitle

\section{Directional limits under exponentially modified prior tails}
\label{supp:sec:exponential-tails}

This section extends the recovery result in the main article to prior
tails containing an exponential factor.  The limiting object is no
longer necessarily the symmetric Lebesgue mixture.  Instead, the
exponential rate in the tail survives as a directional exponential
tilt of Lebesgue measure.  The regularly varying result in the main
article is recovered by setting the exponential rates to zero.

\subsection{Setup and exponentially tilted benchmarks}
\label{supp:subsec:setup}

Let $Y_1,Y_2,\ldots$ be independent observations from
$\mathcal N(\mu,\sigma^2)$, and write
$\theta^\star=\mu/\sigma$ for the standardized effect.  For $n\geq2$,
define
\begin{equation}
\label{supp:eq:Un}
  U_n
  =
  \frac{\sum_{i=1}^nY_i}
       {\sqrt{\sum_{i=1}^nY_i^2}}.
\end{equation}
Whenever the first $n$ observations have positive sample variance,
$|U_n|<\sqrt n$.

Let $R_n\sim\chi_n$ and define
\begin{equation}
\label{supp:eq:hnu}
  A_n(z)=\E[e^{zR_n}],
  \qquad
  h_n(\theta,u)
  =
  e^{-n\theta^2/2}A_n(u\theta),
  \qquad
  u\in(-\sqrt n,\sqrt n).
\end{equation}
Let $\pi$ be a baseline probability density on $\R$ and consider its
location family
\begin{equation}
\label{supp:eq:shifted-prior}
  \pi_m(\theta)=\pi(\theta-m),
  \qquad m\in\R.
\end{equation}
The corresponding scale-invariant likelihood-ratio mixture is
\begin{equation}
\label{supp:eq:shifted-mixture}
  M_{n,m}(u)
  =
  \int_{\R}h_n(\theta,u)\pi(\theta-m)\,\dd\theta,
  \qquad
  M_{n,m}=M_{n,m}(U_n).
\end{equation}
Its extended-Ville calibration and the associated fixed-time and
running $p$-variables are
\begin{align}
  g_m(c)
  &={}
  \E_{0,1}\left[
    \min\left\{1,\frac{M_{2,m}}{c}\right\}
  \right],
  \qquad c>0,
  \label{supp:eq:gm}\\
  P_{n,m}^{E}
  &={}
  g_m(M_{n,m}),
  \qquad
  \widetilde P_{n,m}^{E}
  =
  g_m\left(\max_{2\leq k\leq n}M_{k,m}\right),
  \label{supp:eq:shifted-p-values}
\end{align}
where $g_m(0)=1$ and $g_m(\infty)=0$.

For $\lambda\in\R$, define the $\sigma$-finite mixing measure
\begin{equation}
\label{supp:eq:tilted-measure}
  \Pi_\lambda(\dd\theta)
  =
  e^{\lambda\theta}\,\dd\theta
\end{equation}
and the corresponding mixture
\begin{equation}
\label{supp:eq:tilted-mixture}
  M_n^{(\lambda)}(u)
  =
  \int_{\R}h_n(\theta,u)e^{\lambda\theta}\,\dd\theta,
  \qquad
  M_n^{(\lambda)}=M_n^{(\lambda)}(U_n).
\end{equation}
The following elementary bound ensures that this integral is finite.

\begin{lemma}[Gaussian envelope]
\label{supp:lem:gaussian-envelope}
For every $n\geq2$ and $u\in(-\sqrt n,\sqrt n)$, there exist constants
$C_{n,u},\rho_{n,u}>0$ such that
\begin{equation}
\label{supp:eq:gaussian-envelope}
  0\leq h_n(\theta,u)
  \leq C_{n,u}e^{-\rho_{n,u}\theta^2},
  \qquad \theta\in\R.
\end{equation}
Consequently, $M_n^{(\lambda)}(u)<\infty$ for every fixed
$\lambda\in\R$.
\end{lemma}

\begin{proof}
Using the density of $R_n$ and completing the square gives
\[
  A_n(u\theta)
  =
  c_n e^{u^2\theta^2/2}
  \int_0^\infty
    r^{n-1}e^{-(r-u\theta)^2/2}\,\dd r
\]
for a constant $c_n>0$.  Since
$r^{n-1}\leq C_n\{ |r-u\theta|^{n-1}+|u\theta|^{n-1}\}$,
\[
  A_n(u\theta)
  \leq
  C_{n,u}'(1+|\theta|^{n-1})e^{u^2\theta^2/2}.
\]
It follows that
\[
  h_n(\theta,u)
  \leq
  C_{n,u}'(1+|\theta|^{n-1})
  \exp\left\{-\frac{n-u^2}{2}\theta^2\right\}.
\]
Because $n-u^2>0$, the polynomial factor can be absorbed into a
Gaussian with a smaller exponent, proving
\eqref{supp:eq:gaussian-envelope}.  Multiplication by
$e^{\lambda\theta}$ preserves integrability.
\end{proof}

Define the tilted calibration
\begin{equation}
\label{supp:eq:tilted-calibration}
  g_\lambda(c)
  =
  \E_{0,1}\left[
    \min\left\{1,\frac{M_2^{(\lambda)}}{c}\right\}
  \right],
  \qquad c>0,
\end{equation}
with $g_\lambda(0)=1$ and $g_\lambda(\infty)=0$, and set
\begin{equation}
\label{supp:eq:tilted-p-processes}
  P_n^{(\lambda)}
  =
  g_\lambda\left(M_n^{(\lambda)}\right),
  \qquad
  \widetilde P_n^{(\lambda)}
  =
  g_\lambda\left(
    \max_{2\leq k\leq n}M_k^{(\lambda)}
  \right).
\end{equation}
These are the extended-Ville calibrated $p$-processes associated with
the exponentially tilted mixing measure.  At $\lambda=0$,
\begin{equation}
\label{supp:eq:lambda-zero}
  M_n^{(0)}=M_n^{\Leb},
  \qquad
  g_0=g_{\Leb},
  \qquad
  P_n^{(0)}=P_n^{\Leb}.
\end{equation}

The tilted mixture has an explicit representation.  Tonelli's theorem
and the Gaussian integral give
\begin{equation}
\label{supp:eq:tilted-mixture-expectation}
  M_n^{(\lambda)}(u)
  =
  \sqrt{\frac{2\pi}{n}}
  e^{\lambda^2/(2n)}
  \E\left[
    \exp\left\{
      \frac{u^2R_n^2}{2n}
      +
      \frac{\lambda uR_n}{n}
    \right\}
  \right].
\end{equation}
Equivalently, writing $D_{-n}$ for the parabolic cylinder function,
\begin{align}
  M_n^{(\lambda)}(u)
  &={}
  \sqrt{\frac{2\pi}{n}}
  \left(1-\frac{u^2}{n}\right)^{-n/2}
  \exp\left\{
    \frac{\lambda^2}{2n}
    +
    \frac{\lambda^2u^2}{4n(n-u^2)}
  \right\}
  \nonumber\\
  &\quad\times
  \frac{
    D_{-n}\left(
      -\dfrac{\lambda u}{\sqrt{n(n-u^2)}}
    \right)
  }{
    D_{-n}(0)
  }.
\label{supp:eq:tilted-mixture-closed-form}
\end{align}
In particular,
$M_n^{(-\lambda)}(u)=M_n^{(\lambda)}(-u)$, displaying the
directional nature of the benchmark.

\subsection{Directional recovery theorem}
\label{supp:subsec:directional-recovery}

\begin{assumption}[Prior tails with an exponential factor]
\label{supp:ass:exponential-tails}
There exist $\kappa_+,\kappa_-\geq0$,
$\beta_+,\beta_-\in\R$, and positive slowly varying functions
$L_+$ and $L_-$ such that, as $x\to\infty$,
\begin{equation}
\label{supp:eq:exponential-tail-assumption}
  \pi(x)
  \sim
  e^{-\kappa_+x}x^{-\beta_+}L_+(x),
  \qquad
  \pi(-x)
  \sim
  e^{-\kappa_-x}x^{-\beta_-}L_-(x).
\end{equation}
\end{assumption}

When $\kappa_+=0$ or $\kappa_-=0$, properness of $\pi$ imposes the
corresponding integrability condition on the polynomial and slowly
varying factors.  The regularly varying assumption in the main article
is recovered by taking $\kappa_+=\kappa_-=0$ and
$\beta_+=\beta_-=\beta>1$.

For $s\in\{-1,+1\}$, define the directional exponential rate
\begin{equation}
\label{supp:eq:directional-lambda}
  \lambda_s
  =
  \begin{cases}
    \kappa_-, & s=+1,\\
    -\kappa_+, & s=-1.
  \end{cases}
\end{equation}
Thus $s=+1$ corresponds to $m\to+\infty$, which probes the left tail
of $\pi$, whereas $s=-1$ corresponds to $m\to-\infty$, which probes
the right tail.  Let
\begin{equation}
\label{supp:eq:directional-scale}
  a_m
  =
  \pi_m(0)
  =
  \pi(-m).
\end{equation}
Under Assumption~\ref{supp:ass:exponential-tails},
\begin{equation}
\label{supp:eq:directional-scale-asymptotics}
  a_m
  \sim
  \begin{cases}
    e^{-\kappa_-m}m^{-\beta_-}L_-(m),
      &m\to+\infty,\\[3pt]
    e^{-\kappa_+|m|}|m|^{-\beta_+}L_+(|m|),
      &m\to-\infty.
  \end{cases}
\end{equation}

\begin{theorem}[Directional recovery under exponentially modified tails]
\label{supp:thm:exponential-tail-recovery}
Fix $s\in\{-1,+1\}$ and write $m=sr$, with $r\to\infty$.  Let
$n\geq2$ and fix a dataset $y_{1:n}$ with positive sample variance, so
that
\begin{equation}
\label{supp:eq:fixed-un}
  u_n
  =
  \frac{\sum_{i=1}^ny_i}
       {\sqrt{\sum_{i=1}^ny_i^2}}
  \in(-\sqrt n,\sqrt n).
\end{equation}
Under Assumption~\ref{supp:ass:exponential-tails},
\begin{align}
  \frac{M_{n,sr}}{a_{sr}}
  &\longrightarrow
  M_n^{(\lambda_s)},
  \label{supp:eq:evidence-directional-limit}\\
  g_{sr}(a_{sr}c)
  &\longrightarrow
  g_{\lambda_s}(c),
  \qquad c>0,
  \label{supp:eq:calibration-directional-limit}\\
  P_{n,sr}^{E}
  =
  g_{sr}(M_{n,sr})
  &\longrightarrow
  g_{\lambda_s}\left(M_n^{(\lambda_s)}\right)
  =
  P_n^{(\lambda_s)}.
  \label{supp:eq:p-directional-limit}
\end{align}
If every prefix $y_{1:k}$, $2\leq k\leq n$, has positive sample
variance, then also
\begin{equation}
\label{supp:eq:running-directional-limit}
  \widetilde P_{n,sr}^{E}
  \longrightarrow
  \widetilde P_n^{(\lambda_s)}.
\end{equation}
\end{theorem}

The theorem shows that extended Ville calibration does not remove the
exponential tail rate.  Under $m\to+\infty$, the limiting mixing
measure is $e^{\kappa_-\theta}\dd\theta$; under $m\to-\infty$, it is
$e^{-\kappa_+\theta}\dd\theta$.  The directional benchmark reduces to
Lebesgue measure only when the relevant exponential rate is zero.

\begin{corollary}[Laplace location priors]
\label{supp:cor:laplace}
Let
\begin{equation}
\label{supp:eq:laplace-density}
  \pi(z)
  =
  \frac{1}{2b}e^{-|z|/b},
  \qquad b>0.
\end{equation}
Then
\begin{equation}
\label{supp:eq:laplace-scale}
  a_m
  =
  \frac{1}{2b}e^{-|m|/b},
\end{equation}
and, for every fixed dataset satisfying the conditions of
Theorem~\ref{supp:thm:exponential-tail-recovery},
\begin{align}
  \frac{M_{n,m}}{a_m}
  &\longrightarrow M_n^{(1/b)},
  &
  P_{n,m}^{E}
  &\longrightarrow P_n^{(1/b)},
  &&m\to+\infty,
  \label{supp:eq:laplace-positive-limit}\\
  \frac{M_{n,m}}{a_m}
  &\longrightarrow M_n^{(-1/b)},
  &
  P_{n,m}^{E}
  &\longrightarrow P_n^{(-1/b)},
  &&m\to-\infty.
  \label{supp:eq:laplace-negative-limit}
\end{align}
Thus a Laplace prior has a nondegenerate but directional fallback,
rather than the symmetric Lebesgue fallback obtained for regularly
varying tails.
\end{corollary}

\subsection{Proofs}
\label{supp:subsec:proofs}

\begin{proof}[Proof of Theorem~\ref{supp:thm:exponential-tail-recovery}]
Fix $s\in\{-1,+1\}$.  To simplify notation, let
\[
  (\kappa,\beta,L)
  =
  \begin{cases}
    (\kappa_-,\beta_-,L_-), &s=+1,\\
    (\kappa_+,\beta_+,L_+), &s=-1,
  \end{cases}
  \qquad
  \lambda_s=s\kappa.
\]
For $x>0$, define
\[
  f_s(x)=\pi(-sx),
  \qquad
  q_s(x)=e^{\kappa x}f_s(x).
\]
Assumption~\ref{supp:ass:exponential-tails} gives
\[
  q_s(x)\sim x^{-\beta}L(x),
  \qquad x\to\infty,
\]
so $q_s$ is regularly varying with index $-\beta$.  Moreover,
\begin{equation}
\label{supp:eq:tail-ratio-proof}
  a_{sr}=f_s(r),
  \qquad
  \frac{\pi(\theta-sr)}{a_{sr}}
  =
  \frac{f_s(r-s\theta)}{f_s(r)}
  =
  e^{s\kappa\theta}
  \frac{q_s(r-s\theta)}{q_s(r)}.
\end{equation}
For every fixed $R<\infty$, the uniform convergence theorem for
regularly varying functions \citep[Theorem~1.2.1]{Bingham1987} yields
\begin{equation}
\label{supp:eq:local-exponential-shape}
  \sup_{|\theta|\leq R}
  \left|
    \frac{\pi(\theta-sr)}{a_{sr}}
    -e^{\lambda_s\theta}
  \right|
  \longrightarrow0.
\end{equation}

We first prove the evidence limit.  Fix $R>0$ and, for $r>2R$,
decompose
\begin{align}
  \frac{M_{n,sr}(u)}{a_{sr}}
  &={}
  \int_{|\theta|\leq R}
    h_n(\theta,u)
    \frac{\pi(\theta-sr)}{a_{sr}}\,\dd\theta
  \nonumber\\
  &\quad+
  \int_{R<|\theta|\leq r/2}
    h_n(\theta,u)
    \frac{\pi(\theta-sr)}{a_{sr}}\,\dd\theta
  \nonumber\\
  &\quad+
  \int_{|\theta|>r/2}
    h_n(\theta,u)
    \frac{\pi(\theta-sr)}{a_{sr}}\,\dd\theta
  \nonumber\\
  &=:
  I_{1,s}(R,r)+I_{2,s}(R,r)+I_{3,s}(r).
\label{supp:eq:three-part-decomposition}
\end{align}
By \eqref{supp:eq:local-exponential-shape},
\begin{equation}
\label{supp:eq:I1-limit}
  I_{1,s}(R,r)
  \longrightarrow
  \int_{|\theta|\leq R}
    h_n(\theta,u)e^{\lambda_s\theta}\,\dd\theta.
\end{equation}
On $|\theta|\leq r/2$, the ratio $(r-s\theta)/r$ lies in
$[1/2,3/2]$.  The uniform convergence theorem, applied to $q_s$,
therefore gives a constant $C_\star<\infty$ such that, for all
sufficiently large $r$,
\[
  \frac{q_s(r-s\theta)}{q_s(r)}
  \leq C_\star,
  \qquad |\theta|\leq r/2.
\]
Using \eqref{supp:eq:tail-ratio-proof},
\begin{equation}
\label{supp:eq:I2-bound}
  I_{2,s}(R,r)
  \leq
  C_\star
  \int_{|\theta|>R}
    h_n(\theta,u)e^{\kappa|\theta|}\,\dd\theta.
\end{equation}
Lemma~\ref{supp:lem:gaussian-envelope} makes the last integrand
integrable, and hence the upper bound tends to zero as $R\to\infty$.

For the final term, properness of $\pi$ gives
\begin{equation}
\label{supp:eq:I3-bound}
  I_{3,s}(r)
  \leq
  a_{sr}^{-1}
  \sup_{|\theta|>r/2}h_n(\theta,u).
\end{equation}
By Lemma~\ref{supp:lem:gaussian-envelope}, the supremum is bounded by
$C_{n,u}e^{-\rho_{n,u}r^2/4}$.  Moreover, for every $\varepsilon>0$,
Potter's bound \citep[Theorem~1.5.6]{Bingham1987} and
Assumption~\ref{supp:ass:exponential-tails} imply
\begin{equation}
\label{supp:eq:scale-inverse-bound}
  a_{sr}^{-1}
  =
  O\left(e^{\kappa r}r^{|\beta|+\varepsilon}\right).
\end{equation}
Consequently,
\begin{equation}
\label{supp:eq:I3-limit}
  I_{3,s}(r)
  =
  O\left(
    e^{\kappa r}r^{|\beta|+\varepsilon}
    e^{-\rho_{n,u}r^2/4}
  \right)
  \longrightarrow0.
\end{equation}
Combining \eqref{supp:eq:I1-limit}--\eqref{supp:eq:I3-limit} and then
letting $R\to\infty$ proves
\[
  \frac{M_{n,sr}(u)}{a_{sr}}
  \longrightarrow
  \int_{\R}h_n(\theta,u)e^{\lambda_s\theta}\,\dd\theta
  =
  M_n^{(\lambda_s)}(u),
\]
which is \eqref{supp:eq:evidence-directional-limit}.

Under the null, $|U_2|<\sqrt2$ almost surely.  Applying the preceding
limit with $n=2$ and $u=U_2$ gives
\[
  \frac{M_{2,sr}}{a_{sr}}
  \longrightarrow
  M_2^{(\lambda_s)}
  \qquad
  \Prob_{0,1}\text{-almost surely}.
\]
For every $c>0$, bounded convergence therefore gives
\begin{align}
  g_{sr}(a_{sr}c)
  &={}
  \E_{0,1}\left[
    \min\left\{
      1,
      \frac{M_{2,sr}}{a_{sr}c}
    \right\}
  \right]
  \nonumber\\
  &\longrightarrow
  \E_{0,1}\left[
    \min\left\{
      1,
      \frac{M_2^{(\lambda_s)}}{c}
    \right\}
  \right]
  =
  g_{\lambda_s}(c),
\label{supp:eq:calibration-limit-proof}
\end{align}
proving \eqref{supp:eq:calibration-directional-limit}.

For the fixed dataset, put
\[
  c_r=\frac{M_{n,sr}}{a_{sr}},
  \qquad
  c=M_n^{(\lambda_s)}.
\]
Then $c_r\to c\in(0,\infty)$.  Since every $g_{sr}$ is
non-increasing, for $\varepsilon\in(0,c)$ and all sufficiently large
$r$,
\[
  g_{sr}\left(a_{sr}(c+\varepsilon)\right)
  \leq
  g_{sr}(a_{sr}c_r)
  \leq
  g_{sr}\left(a_{sr}(c-\varepsilon)\right).
\]
Applying \eqref{supp:eq:calibration-directional-limit} to the outer
terms and then letting $\varepsilon\downarrow0$ gives
\[
  g_{sr}(M_{n,sr})
  \longrightarrow
  g_{\lambda_s}\left(M_n^{(\lambda_s)}\right),
\]
because $g_{\lambda_s}$ is continuous on $(0,\infty)$.  This proves
\eqref{supp:eq:p-directional-limit}.  If every prefix has positive
sample variance, the evidence limit holds for each $k=2,\ldots,n$.
Taking the maximum over this finite set and repeating the same
monotonicity argument proves
\eqref{supp:eq:running-directional-limit}.
\end{proof}

\begin{proof}[Derivation of
\eqref{supp:eq:tilted-mixture-expectation} and
\eqref{supp:eq:tilted-mixture-closed-form}]
Tonelli's theorem and the Gaussian integral give
\begin{align}
  M_n^{(\lambda)}(u)
  &={}
  \E\left[
    \int_{\R}
    \exp\left\{
      -\frac{n\theta^2}{2}
      +(uR_n+\lambda)\theta
    \right\}\,\dd\theta
  \right]
  \nonumber\\
  &={}
  \sqrt{\frac{2\pi}{n}}
  \E\left[
    \exp\left\{
      \frac{(uR_n+\lambda)^2}{2n}
    \right\}
  \right],
\label{supp:eq:gaussian-integral-derivation}
\end{align}
which is \eqref{supp:eq:tilted-mixture-expectation}.  To obtain
\eqref{supp:eq:tilted-mixture-closed-form}, use the density of $R_n$
and the identity
\begin{equation}
\label{supp:eq:parabolic-cylinder-identity}
  \int_0^\infty
  t^{\nu-1}e^{-t^2/2-zt}\,\dd t
  =
  \Gamma(\nu)e^{z^2/4}D_{-\nu}(z),
  \qquad \nu>0.
\end{equation}
After the change of variables
$t=\sqrt{1-u^2/n}\,R_n$, this gives the stated formula.  The
normalizing identity
\begin{equation}
\label{supp:eq:parabolic-cylinder-zero}
  D_{-n}(0)
  =
  \frac{2^{-n/2}\sqrt\pi}{\Gamma((n+1)/2)}
\end{equation}
and Legendre's duplication formula ensure that the expression reduces
to $M_n^{\Leb}(u)$ when $\lambda=0$.
\end{proof}

\bibliographystyle{abbrvnat}
\bibliography{robustttest}

\end{document}
